\section{The irrationality contradiction}\label{sec:contradiction}

We now prove Proposition~\ref{prop:criterion}. Assume its two
inequalities hold and suppose that at most $B-1$ of
$\beta(m+1),\beta(m+3),\ldots,\beta(S-1)$ are irrational.
Choose $B-1$ distinct even indices
$\nu_1<\cdots<\nu_{B-1}$ containing all possible irrational
indices, adding further indices if necessary. There are enough
choices because $S>mDL$ and $D\ge B$.

The coefficients $c_j$ have already fixed the common multiplier
$W_n$ and the seed vector. As in the elimination argument
of~\cite[Section~5]{LY20}, we choose fixed integer weights
$w_1,\ldots,w_B$ on the $B$ forms to annihilate the columns
\[
 (1^{\nu_j},2^{\nu_j},\ldots,B^{\nu_j})^{\mathsf T},
                                  \qquad 1\le j\le B-1.
\]
Use their signed maximal minors. The last weight $w_B$ is nonzero:
its minor is a generalized Vandermonde determinant at
$1^2,\ldots,(B-1)^2$ with exponents $\nu_j/2$.
A vanishing minor would give a nonzero polynomial with at most
$B-1$ monomials and $B-1$ distinct positive roots. A polynomial
with $q$ monomials has at most $q-1$ such roots: factor out its
smallest power and apply Rolle's theorem inductively to its
derivative. This proves nonsingularity. For $B=1$, take $w_1=1$.

Set
\[
 H_n=\sum_{b=1}^B w_bF_{n,b},\qquad
 q_w(n)=\sum_{b=1}^B w_bq_b(n).
\]
The chosen beta coefficients cancel at every index. All remaining
beta values are rational under the assumption, so $H_n$ is a
rational sequence. Let $D_0$ clear the denominators of those
rational beta values. Proposition~\ref{prop:denominators} and
\eqref{eq:cleared-decay} give
\begin{equation}\label{eq:integers}
 D_0\delta_n^{S+k}H_n\in\Z,\qquad
 D_0\delta_n^{S+k}H_n\longrightarrow0
                   \quad(n\to\infty,\ n\text{ odd}).
\end{equation}

The seed choice \eqref{eq:seed-target} gives
$q_w(1)=w_B\eta\ne0$. Exclude the finitely many primes dividing
$D_0$ or the numerator or denominator of $q_w(1)$, and take $p$
large enough for Proposition~\ref{prop:dilation}. Its constant
congruence at $\ell=1$ gives
\[
 p^\tau q_w(p)\equiv\ch(p)2^{-S}q_w(1)\not\equiv0\pmod p.
\]
Every remaining beta term of $p^\tau H_p$ is divisible by $p$
in $\Zp$, by \eqref{eq:rho-dilation} and the assumed rationality.
Thus $H_p\ne0$ for infinitely many primes. Equation~\eqref{eq:integers}
gives infinitely many nonzero integers tending to zero, which is
impossible. This proves Proposition~\ref{prop:criterion}.
